% GENERATED FILE. Edit canonical lessons, then run npm run generate:books.
% canonical-source-hash: fnv1a64:5747222c3bd63524
% canonical-lessons: TE-C522-vrtti, TE-C522-hoda, TE-C522-padavi, TE-C522-padonnati, TE-C522-asayam

\chapter{A Profession}
\label{ch:te-a-profession}
\hlchaptermodality{522}

\begin{chapteropening}
\textbf{By the end of this chapter:} \emph{I can talk about my profession, my position and my ambition.}

Complete the last lesson of chapter 522: I can talk about my profession, my position and my ambition.
\end{chapteropening}

\section[\texorpdfstring{\te{వృత్తి} (vṛtti)}{వృత్తి (vṛtti)}]{\te{వృత్తి} (vṛtti) — a profession; a calling}

\label{lesson:TE-C522-vrtti}



\begin{quote}
Before the new one: say the Telugu for to put on weight, then the Telugu for nutrition.
\end{quote}

\subsection*{You'll want to know: \te{వృత్తి}}
\textbf{\te{వృత్తి}} — \emph{vṛtti} — \textquotedblleft{}a profession; a calling\textquotedblright{}.

\begin{quote}
\te{మా} \te{నాన్నగారి} \te{వృత్తి} \te{వైద్యం}.
\end{quote}

\begin{quote}
\emph{mā nānnagāri vṛtti vaidyaṁ.}
\end{quote}

\begin{quote}
My father's profession is medicine.
\end{quote}

Your calling: \te{వృత్తిరీత్యా}, by profession.

\begin{cousinweb}
\textbf{\te{వృత్తి}} (\emph{vṛtti}) is borrowed from Sanskrit \emph{vṛtti}, \textquotedblleft{}mode of life, course of action\textquotedblright{} (Monier-Williams): the way you make your living.
\end{cousinweb}

\begin{grammarlens}[title={What you've built}]
One more word for describing what you lived and felt.
\end{grammarlens}

\subsection*{Your turn}
\begin{itemize}
\raggedright
  \item \emph{Say it:} \emph{vṛtti}
  \item \emph{Say it:} \emph{vṛtti}, once more
  \item \emph{Say it:} \emph{mā nānnagāri vṛtti vaidyaṁ.} — My father's profession is medicine.
  \item \emph{Recall:} say the Telugu for to put on weight, then the Telugu for nutrition, then say \emph{vṛtti} again
\end{itemize}

\begin{tcolorbox}[breakable,colback=teal!4,colframe=teal!35!black,title={Before you move on}]
What does \te{వృత్తి} mean? (\textquotedblleft{}A profession; a calling\textquotedblright{}.) Say it once more.
\end{tcolorbox}
\section[\texorpdfstring{\te{హోదా} (hōdā)}{హోదా (hōdā)}]{\te{హోదా} (hōdā) — status; a rank}

\label{lesson:TE-C522-hoda}



\begin{quote}
Before the new one: say the Telugu for nutrition, then the Telugu for a profession.
\end{quote}

\subsection*{You'll want to know: \te{హోదా}}
\textbf{\te{హోదా}} — \emph{hōdā} — \textquotedblleft{}status; a rank\textquotedblright{}.

\begin{quote}
\te{కొత్త} \te{ఉద్యోగంలో} \te{అతని} \te{హోదా} \te{పెరిగింది}.
\end{quote}

\begin{quote}
\emph{kotta udyōgaṁlō atani hōdā perigindi.}
\end{quote}

\begin{quote}
In his new job his status rose.
\end{quote}

Rank in an office or in society.

\begin{cousinweb}
\textbf{\te{హోదా}} (\emph{hōdā}) came into Telugu through Hindustani \emph{ohdā}, \textquotedblleft{}an office, a rank\textquotedblright{}, which is usually traced to Arabic.
\end{cousinweb}

\begin{grammarlens}[title={What you've built}]
One more word for describing what you lived and felt.
\end{grammarlens}

\subsection*{Your turn}
\begin{itemize}
\raggedright
  \item \emph{Say it:} \emph{hōdā}
  \item \emph{Say it:} \emph{hōdā}, once more
  \item \emph{Say it:} \emph{kotta udyōgaṁlō atani hōdā perigindi.} — In his new job his status rose.
  \item \emph{Recall:} say the Telugu for nutrition, then the Telugu for a profession, then say \emph{hōdā} again
\end{itemize}

\begin{tcolorbox}[breakable,colback=teal!4,colframe=teal!35!black,title={Before you move on}]
What does \te{హోదా} mean? (\textquotedblleft{}Status; a rank\textquotedblright{}.) Say it once more.
\end{tcolorbox}
\section[\texorpdfstring{\te{పదవి} (padavi)}{పదవి (padavi)}]{\te{పదవి} (padavi) — a position; a post}

\label{lesson:TE-C522-padavi}



\begin{quote}
Before the new one: say the Telugu for a profession, then the Telugu for status.
\end{quote}

\subsection*{You'll want to know: \te{పదవి}}
\textbf{\te{పదవి}} — \emph{padavi} — \textquotedblleft{}a position; a post\textquotedblright{}.

\begin{quote}
\te{ఆమె} \te{కంపెనీలో} \te{పెద్ద} \te{పదవిలో} \te{ఉంది}.
\end{quote}

\begin{quote}
\emph{āme kampenīlō pedda padavilō undi.}
\end{quote}

\begin{quote}
She holds a senior position in the company.
\end{quote}

\te{పదవిలో} \te{ఉండు}, to hold a post.

\begin{cousinweb}
\textbf{\te{పదవి}} (\emph{padavi}) is borrowed from Sanskrit \emph{padavī}, \textquotedblleft{}a road, a track; a station, a place\textquotedblright{} (Monier-Williams).
\end{cousinweb}

\begin{grammarlens}[title={What you've built}]
One more word for describing what you lived and felt.
\end{grammarlens}

\subsection*{Your turn}
\begin{itemize}
\raggedright
  \item \emph{Say it:} \emph{padavi}
  \item \emph{Say it:} \emph{padavi}, once more
  \item \emph{Say it:} \emph{āme kampenīlō pedda padavilō undi.} — She holds a senior position in the company.
  \item \emph{Recall:} say the Telugu for a profession, then the Telugu for status, then say \emph{padavi} again
\end{itemize}

\begin{tcolorbox}[breakable,colback=teal!4,colframe=teal!35!black,title={Before you move on}]
What does \te{పదవి} mean? (\textquotedblleft{}A position; a post\textquotedblright{}.) Say it once more.
\end{tcolorbox}
\section[\texorpdfstring{\te{పదోన్నతి} (padōnnati)}{పదోన్నతి (padōnnati)}]{\te{పదోన్నతి} (padōnnati) — a promotion}

\label{lesson:TE-C522-padonnati}



\begin{quote}
Before the new one: say the Telugu for status, then the Telugu for a position.
\end{quote}

\subsection*{You'll want to know: \te{పదోన్నతి}}
\textbf{\te{పదోన్నతి}} — \emph{padōnnati} — \textquotedblleft{}a promotion\textquotedblright{}.

\begin{quote}
\te{కష్టపడి} \te{పనిచేసినందుకు} \te{ఆమెకు} \te{పదోన్నతి} \te{వచ్చింది}.
\end{quote}

\begin{quote}
\emph{kaṣṭapaḍi panicēsinanduku āmeku padōnnati vaccindi.}
\end{quote}

\begin{quote}
She got a promotion for working hard.
\end{quote}

Formal; in speech many say \te{ప్రమోషను}.

\begin{cousinweb}
\textbf{\te{పదోన్నతి}} (\emph{padōnnati}) is a learned compound built from Sanskrit \emph{pada}, \textquotedblleft{}step, position\textquotedblright{}, and \emph{unnati}, \textquotedblleft{}rising, advancement\textquotedblright{} (Monier-Williams).
\end{cousinweb}

\begin{grammarlens}[title={What you've built}]
One more word for describing what you lived and felt.
\end{grammarlens}

\subsection*{Your turn}
\begin{itemize}
\raggedright
  \item \emph{Say it:} \emph{padōnnati}
  \item \emph{Say it:} \emph{padōnnati}, once more
  \item \emph{Say it:} \emph{kaṣṭapaḍi panicēsinanduku āmeku padōnnati vaccindi.} — She got a promotion for working hard.
  \item \emph{Recall:} say the Telugu for status, then the Telugu for a position, then say \emph{padōnnati} again
\end{itemize}

\begin{tcolorbox}[breakable,colback=teal!4,colframe=teal!35!black,title={Before you move on}]
What does \te{పదోన్నతి} mean? (\textquotedblleft{}A promotion\textquotedblright{}.) Say it once more.
\end{tcolorbox}
\section[\texorpdfstring{\te{ఆశయం} (āśayaṁ)}{ఆశయం (āśayaṁ)}]{\te{ఆశయం} (āśayaṁ) — an ambition; an ideal}

\label{lesson:TE-C522-asayam}



\begin{quote}
Before the new one: say the Telugu for a position, then the Telugu for a promotion.
\end{quote}

\subsection*{You'll want to know: \te{ఆశయం}}
\textbf{\te{ఆశయం}} — \emph{āśayaṁ} — \textquotedblleft{}an ambition; an ideal\textquotedblright{}.

\begin{quote}
\te{డాక్టరు} \te{కావడమే} \te{నా} \te{ఆశయం}.
\end{quote}

\begin{quote}
\emph{ḍākṭaru kāvaḍamē nā āśayaṁ.}
\end{quote}

\begin{quote}
Becoming a doctor is my ambition.
\end{quote}

A life goal, nobler than a passing wish.

\begin{cousinweb}
\textbf{\te{ఆశయం}} (\emph{āśayaṁ}) is borrowed from Sanskrit \emph{āśaya} (CDIAL 1453), which first meant a resting place; in Telugu it names the aim a person lives for.
\end{cousinweb}

\begin{grammarlens}[title={What you've built}]
One more word for describing what you lived and felt.
\end{grammarlens}

\subsection*{Your turn}
\begin{itemize}
\raggedright
  \item \emph{Say it:} \emph{āśayaṁ}
  \item \emph{Say it:} \emph{āśayaṁ}, once more
  \item \emph{Say it:} \emph{ḍākṭaru kāvaḍamē nā āśayaṁ.} — Becoming a doctor is my ambition.
  \item \emph{Recall:} say the Telugu for a position, then the Telugu for a promotion, then say \emph{āśayaṁ} again
\end{itemize}

\begin{tcolorbox}[breakable,colback=teal!4,colframe=teal!35!black,title={Before you move on}]
What does \te{ఆశయం} mean? (\textquotedblleft{}An ambition; an ideal\textquotedblright{}.) Say it once more.
\end{tcolorbox}
